%% \PassOptionsToPackage{hidelinks,unicode}{hyperref} % Example options

%%%%%%%%% PACKAGES

\usepackage{booktabs}     % for nice horizontal lines
%\usepackage{caption}

% --- Preamble (place in your preamble) ---
% \usepackage[switch]{lineno} 
\usepackage{graphicx}
\usepackage{subcaption}
\usepackage{placeins} % for \FloatBarrier to control float placement
\usepackage{float}     % for [H] exact float placement
% Independent "Listing N" float with a two-column (starred) variant for the
% motivating-example diff. newfloat only needs keyval, so it does not pull in
% the caption package (which conflicts with acmart).
\usepackage{newfloat}
\DeclareFloatingEnvironment[name=Listing, fileext=lol, placement=tbp, within=none]{listing}
\usepackage{tcolorbox}
\tcbuselibrary{skins,breakable,listings}
\usepackage{listings}
\usepackage{xcolor}  % table option passed via \PassOptionsToPackage in main.tex
\usepackage{array}
% Listings styles
\lstdefinestyle{code}{
  basicstyle=\ttfamily\small,
  numbers=left, numberstyle=\scriptsize, numbersep=6pt,
  breaklines=true, breakatwhitespace=true,
  frame=single, framerule=0.4pt,
  tabsize=2, showstringspaces=false
}
\lstdefinestyle{cpp}{style=code, language=C++}
\lstdefinestyle{json}{style=code, language=}
\lstdefinestyle{txt}{style=code, language=}

% Ensure circled step numbers in listings render reliably
% Maps Unicode circled digits to the TikZ-based \circled{n} macro
\lstset{literate=
  {❶}{{\circled{1}}}1
  {❷}{{\circled{2}}}1
  {❸}{{\circled{3}}}1
  {❹}{{\circled{4}}}1
  {❺}{{\circled{5}}}1
  {❻}{{\circled{6}}}1
  {❼}{{\circled{7}}}1
}

% A reusable “code card” environment
\newtcblisting{codecard}[2][]{%
  enhanced, breakable,
  colback=white, colframe=black!12,
  boxrule=0.6pt, arc=2mm, left=1mm, right=1mm, top=1mm, bottom=1mm,
  title={#2}, fonttitle=\bfseries\footnotesize,
  listing only, listing options={#1}
}

% --- GitHub-style split-diff support (used in the Motivating Example) ---------
% Per-cell background highlighting via colortbl's \cellcolor (loaded explicitly
% here: the class loads xcolor *without* the table option, so \cellcolor/\rowcolor
% are otherwise undefined). Code is typeset with \lstinline so quotes/brackets
% stay verbatim-safe.
\usepackage{colortbl}
\definecolor{diffadd}{RGB}{225,251,231}   % GitHub-ish light green (added)
\definecolor{diffdel}{RGB}{255,233,232}   % GitHub-ish light red   (removed)
% Minimal Swift language (listings has no built-in Swift definition)
\lstdefinelanguage{Swift}{%
  morekeywords={func,let,var,while,for,in,if,else,guard,return,true,false,nil},
  sensitive=true, morestring=[b]", morecomment=[l]{//}}
% Compact inline style for the diff cells; long UI-locator lines need it small.
\lstdefinestyle{swiftmini}{%
  basicstyle=\ttfamily\scriptsize, columns=fullflexible, keepspaces=true,
  language=Swift, keywordstyle=\bfseries\color{blue!60!black},
  literate={"}{{\textquotedbl}}1}
% +/- gutter glyphs and the gray line-number cell
\newcommand{\gplus}{\textcolor{green!55!black}{\ttfamily\scriptsize +}}
\newcommand{\gminus}{\textcolor{red!70!black}{\ttfamily\scriptsize $-$}}
\newcommand{\dnum}[1]{\textcolor{black!40}{\ttfamily\scriptsize #1}}
% Shorthand: shade the three cells of one diff half-row (color, then lineno/gutter/code)
\newcommand{\dadd}{\cellcolor{diffadd}}
\newcommand{\ddel}{\cellcolor{diffdel}}
% Code cell + keyword: \lstinline cannot be used here because colortbl/array
% re-scans cells and breaks its verbatim catcodes, so we typeset code as plain
% monospace (T1 renders ", [, ], <, > correctly; ~ gives visible indentation).
\newcommand{\cd}[1]{{\ttfamily\scriptsize\strut #1}}
\newcommand{\kw}[1]{\textcolor{blue!60!black}{\textbf{#1}}}
% One full diff line in a single (optionally shaded) cell, so the line number,
% the +/- gutter, and the code share one continuous background. Fixed-width
% \makebox keeps the columns aligned and guarantees the gap after the number.
% #1=row color, #2=line number, #3=gutter glyph (\gplus/\gminus or empty), #4=code
\newcommand{\diffline}[4]{%
  \cellcolor{#1}\makebox[1.3em][r]{\dnum{#2}}\hspace{8pt}%
  \makebox[0.6em][c]{#3}\hspace{6pt}\cd{#4}}



%% \usepackage{soul}
% xcolor is loaded by the class with the 'table' option (now explicitly loaded above)
\usepackage{tabularx}
%% %% Custom packages end here

%% \usepackage{cite}
%% \usepackage{amsmath}
%% \usepackage{booktabs}
\usepackage{wrapfig}
\usepackage{multirow}
\usepackage{multicol}
\usepackage{enumitem}
%% \usepackage{color}
\usepackage{listings}
%% \usepackage{graphicx}
%% \usepackage{rotate}
%% %% \usepackage{url}
%% %\usepackage[hidelinks]{hyperref}
%% \usepackage{breakurl}
%% \usepackage{flushend}
\usepackage{xspace}
%% \usepackage[round-mode=places,round-precision=1,group-separator={,},group-minimum-digits={3},output-decimal-marker={.}, round-pad=false, exponent-mode = threshold, exponent-thresholds = -3:5, tight-spacing=true]{siunitx}
\usepackage{subcaption}
%\usepackage{pgf-pie}
%% \usepackage{bbding}
\usepackage{pifont}
%% \usepackage{seqsplit}
%% \usepackage{pgfplots}
%% \pgfplotsset{width=7cm,compat=1.8}
%% \usepackage{datetime2}

%% % Tikz Figs
\usepackage{tikz}
\tikzset{>=latex}
\usetikzlibrary{shapes.geometric,arrows.meta,positioning,fit,backgrounds,calc} % for flow diagrams and icon pics

%%%%%%%%%
% ---------- Icon pics (line-art, ~0.5 cm) ----------
\tikzset{
  ic/.style={line width=0.6pt, line cap=round, line join=round},
  pics/play/.style={code={
    \draw[ic] (0,0) circle (0.26);
    \fill (-0.08,-0.13) -- (-0.08,0.13) -- (0.16,0) -- cycle;
  }},
  pics/search/.style={code={
    \draw[ic] (0.03,0.06) circle (0.16);
    \draw[ic] (-0.09,-0.06) -- (-0.22,-0.19);
  }},
  pics/bug/.style={code={
    \draw[ic] (0,0) ellipse (0.15 and 0.21);
    \draw[ic] (0,0.21) -- (0,0.30);
    \draw[ic] (0,0.30) circle (0.07);
    \draw[ic] (-0.04,0.36) -- (-0.09,0.43);
    \draw[ic] (0.04,0.36) -- (0.09,0.43);
    \foreach \y in {0.10,0,-0.10}{
      \draw[ic] (-0.13,\y) -- (-0.26,\y+0.05);
      \draw[ic] (0.13,\y) -- (0.26,\y+0.05);
    }
  }},
  pics/shield/.style={code={
    \draw[ic] (0,0.24) -- (-0.17,0.17) -- (-0.17,-0.07)
      .. controls (-0.17,-0.20) and (-0.07,-0.28) .. (0,-0.30)
      .. controls (0.07,-0.28) and (0.17,-0.20) .. (0.17,-0.07)
      -- (0.17,0.17) -- cycle;
    \draw[ic] (-0.08,0.0) -- (-0.02,-0.08) -- (0.10,0.10);
  }},
  pics/funnel/.style={code={
    \draw[ic] (-0.22,0.20) -- (0.22,0.20) -- (0.05,-0.05) -- (0.05,-0.22)
      -- (-0.05,-0.22) -- (-0.05,-0.05) -- cycle;
  }},
  pics/db/.style={code={
    \draw[ic] (-0.16,0.18) arc (180:360:0.16 and 0.07);
    \draw[ic] (-0.16,0.18) arc (180:0:0.16 and 0.07);
    \draw[ic] (-0.16,0.18) -- (-0.16,-0.18);
    \draw[ic] (0.16,0.18) -- (0.16,-0.18);
    \draw[ic] (-0.16,0.0) arc (180:360:0.16 and 0.07);
    \draw[ic] (-0.16,-0.18) arc (180:360:0.16 and 0.07);
  }},
  pics/brain/.style={code={
    \draw[ic] (0,0.26)
      .. controls (-0.22,0.26) and (-0.28,0.08) .. (-0.20,-0.02)
      .. controls (-0.30,-0.10) and (-0.20,-0.28) .. (-0.05,-0.22)
      .. controls (0.0,-0.30) and (0.0,-0.30) .. (0.05,-0.22)
      .. controls (0.20,-0.28) and (0.30,-0.10) .. (0.20,-0.02)
      .. controls (0.28,0.08) and (0.22,0.26) .. (0,0.26) -- cycle;
    \draw[ic] (0,0.26) -- (0,-0.24);
    \draw[ic] (-0.10,0.12) .. controls (-0.02,0.10) and (-0.02,0.0) .. (-0.10,-0.04);
    \draw[ic] (0.10,0.12) .. controls (0.02,0.10) and (0.02,0.0) .. (0.10,-0.04);
  }},
  pics/book/.style={code={
    \draw[ic] (0,-0.16) .. controls (-0.10,-0.24) and (-0.24,-0.20) .. (-0.26,-0.16)
      -- (-0.26,0.16) .. controls (-0.20,0.12) and (-0.06,0.10) .. (0,0.18);
    \draw[ic] (0,-0.16) .. controls (0.10,-0.24) and (0.24,-0.20) .. (0.26,-0.16)
      -- (0.26,0.16) .. controls (0.20,0.12) and (0.06,0.10) .. (0,0.18);
    \draw[ic] (0,0.18) -- (0,-0.16);
  }},
  pics/graph/.style={code={
    \draw[ic] (-0.30,0) -- (0.30,0);
    \draw[ic] (0,0) -- (0,0.20);
    \foreach \x in {-0.30,-0.10,0.10,0.30}{ \fill (\x,0) circle (0.035); \draw[ic] (\x,0) circle (0.035);}
    \fill (0,0.20) circle (0.035); \draw[ic] (0,0.20) circle (0.035);
  }},
  pics/repo/.style={code={
    \draw[ic] (-0.10,0.26) -- (-0.10,-0.26);
    \foreach \y in {0.20,-0.04,-0.20}{ \fill (-0.10,\y) circle (0.05); \draw[ic](-0.10,\y) circle (0.05);}
    \draw[ic] (-0.10,-0.04) .. controls (0.06,-0.04) and (0.16,0.04) .. (0.16,0.16);
    \fill (0.16,0.20) circle (0.05); \draw[ic] (0.16,0.20) circle (0.05);
  }},
}
%%%%%%%%%


%% %\usepackage{pgf-umlsd} % for flow diagram
%% \usetikzlibrary{calc}
%% \usetikzlibrary{shapes.geometric}
%% \usetikzlibrary{decorations.pathreplacing}
%% \usetikzlibrary{positioning}
\newcommand*\circled[1]{\tikz[baseline=(char.base)]{
            \node[shape=circle,draw,fill=black,text=white,inner
            sep=0.4pt] (char) {\small #1};}}

%% %%%%%%%%%% CHECKMARKS


%% %%%%%%%% CONCLUDING BOXES
\usepackage[framemethod=tikz]{mdframed}
\mdfdefinestyle{mpdframe}{
   frametitlebackgroundcolor   =black!15,
    frametitlerule              =true,
    roundcorner                 =3pt,
    middlelinewidth             =1pt,
    skipabove                   =\topskip,
    skipbelow                   =\topskip,
    innermargin                 =0.1cm, % smaller makes box bigger
    outermargin                 =0.1cm,
    innerleftmargin             =0.1cm,
    innerrightmargin            =0cm,
    innertopmargin              =0.1cm,
    innerbottommargin           =0.1cm,
    align=center   
}

%%%%%%%%%% CODE and MACROS

%\input{defs-code}
%% \input{defs-logic}

%%%%%%%%%% UNCLASSIFIED
\newcommand{\Caption}[1]{\caption{#1}}
\newcommand{\NAtname}{-}
\newcommand{\specialcell}[2][c]{\begin{tabular}[#1]{@{}c@{}}#2\end{tabular}}
% Right-aligned special cell for numbers with inline deltas
\newcommand{\specialcellr}[1]{\begin{tabular}[c]{@{}r@{}}#1\end{tabular}}
\newcommand{\DefMacro}[2]{%
   \expandafter\newcommand\csname rmk-#1\endcsname{#2}%
}
\newcommand{\UseMacro}[1]{\csname rmk-#1\endcsname}


\newcommand{\YC}{\textcolor{blue}{\textbf{+}}} %{$\checkmark$}}
\newcommand{\NC}{\textcolor{red}{\textbf{-}}} %\textbf{$\times$}}}
\newcommand{\MC}{\textcolor{magenta}{\textbf{$\pm$}}}

\newcommand{\etal}{et al.}
\newcommand{\eg}{e.g.}
\newcommand{\ie}{i.e.}
\newcommand{\Fix}[1]{\textcolor{red}{[#1]}}
\newcommand{\Fixsd}[1]{\textcolor{purple}{SD:[#1]}}
\newcommand{\Bren}[1]{\textcolor{teal}{Breno:~[#1]}}
\newcommand{\Den}[1]{\textcolor{blue}{Denini:~[#1]}}


\newcommand{\TODO}[1]{{\color{red}{\bf TODO:} #1}}
\newcommand{\Comment}[1]{}
\newcommand{\Boldface}[1]{\textbf{#1}}
\newcommand{\Space}[1]{}
\newcommand{\hand}[1]{\textcolor{green}{[#1]}}
\newcommand{\raw}[1]{\textcolor{red}{#1}}
\newcommand{\edit}[1]{#1}

\newcommand{\MeanStd}[2]{#1$\pm$#2}
\newcommand{\MeanStdSec}[2]{#1$\pm$#2 seconds}
\newcommand{\plus}{$+$}

\newcommand{\MyPara}[1]{\vspace{1pt}\noindent\textbf{#1}.}
\newcommand{\MyParaOnly}[1]{\noindent\textbf{#1}}
\newcommand{\MyItPara}[1]{\noindent\textit{#1}.}
\newcommand{\reducedstrut}{\vrule width 0pt height .9\ht\strutbox depth .9\dp\strutbox\relax}
%\newcommand{\InputWithSpace}[1]{\bgroup\def\arraystretch{1.1}\input{#1}\egroup}
\newcommand{\Code}[1]{{\small\ifmmode{\texttt{#1}}\else$\texttt{#1}$\fi}}
\newcommand{\CodeIn}[1]{{\small\ifmmode{\mathtt{#1}}\else$\mathtt{#1}$\fi}}
\newcommand{\ColorBack}[1]{%
  \begingroup \setlength{\fboxsep}{0pt}% \colorbox{purple!20}{\reducedstrut#1\/}% \endgroup
}

%% % ----------
%% % THIS BELOW COPIED FROM vkuncak/doc/vmcai09/defs.tex
\definecolor{gray}{RGB}{211,211,211}



% Define the TODO command with highlight and margin note
\newcommand{\todo}[2][normal]{%
  \ifdefined\showtodos
    \ifstrequal{#1}{high}{\sethlcolor{red}}{}%
    \ifstrequal{#1}{normal}{\sethlcolor{yellow}}{}%
    \ifstrequal{#1}{low}{\sethlcolor{green}}{}%
    {\hl{\textbf{TODO:} #2}}%
    \marginpar{\textcolor{red}{\textbf{TODO}}}%
  \fi
}

% Control visibility of TODOs with this toggle
\newif\ifshowtodos
\showtodostrue  % Set to \showtodosfalse before final submission



%%%%%%

%%
%% end of the preamble, start of the body of the document source.
\setkeys{Gin}{draft=false}
\renewcommand{\lstlistingname}{Listing}
\usepackage{tcolorbox}
\usepackage{multirow}
\usepackage{placeins}
\usepackage[hidelinks]{hyperref}
\usepackage{makecell}
\usepackage{tabularx}
\usepackage{array}
\usepackage{fvextra}
\usepackage{microtype}
\emergencystretch=2em
\raggedbottom

\newenvironment{packed_itemize}{
\vspace{-0.1ex}
\begin{list}{\labelitemi}{\leftmargin=1.7em}
 \setlength{\itemsep}{0pt} \setlength{\parskip}{0pt} \setlength{\parsep}{0pt} \setlength{\headsep}{0pt} \setlength{\topskip}{0pt} \setlength{\topmargin}{0pt} \setlength{\topsep}{0pt} \setlength{\partopsep}{0pt}
 }{\end{list}
\vspace{-0.1ex}
}


\newcommand{\tool}{\texttt{FLAG}}
\newcommand{\pcount}{\textsc{2}}
\newcommand{\x}{x\xspace}
\newcommand{\Contrib}[1]{$\star$#1}
\newcommand{\sota}{SoTA\xspace}

\newcommand{\repository}{https://github.com}

\usepackage{pifont}% http://ctan.org/pkg/pifont
\newcommand{\cmark}{\ding{51}}%
\newcommand{\xmark}{\ding{55}}%

\definecolor{codebg}{rgb}{0.95, 0.95, 0.92} % Light background
\definecolor{keywordcolor}{rgb}{0.0, 0.0, 0.6} % Blue for keywords
\definecolor{stringcolor}{rgb}{0.58, 0.0, 0.82} % Purple for strings
\definecolor{commentcolor}{rgb}{0.0, 0.5, 0.0} % Green for comments 

\lstset{
    language=Go,              % Language
    backgroundcolor=\color{codebg}, % Background color
    basicstyle=\ttfamily\small, % Font type and size
    keywordstyle=\color{keywordcolor}\bfseries, % Style for keywords
    stringstyle=\color{stringcolor}, % Style for strings
    commentstyle=\color{commentcolor}\itshape, % Style for comments
    numbers=left,               % Line numbers on the left
    numberstyle=\tiny,          % Line number font size
    stepnumber=1,               % Line numbering step
    frame=single,               % Frame around the code
    breaklines=true,            % Line breaking
    tabsize=2,                  % Tab size
    showspaces=false,           % Don't display spaces
    showstringspaces=false,     % Don't display string spaces
    captionpos=b,               % Caption position (b for bottom)
    lineskip=1pt,
    linewidth=1\columnwidth,
    xleftmargin=13pt,    % Add padding to the left
    numbersep=5pt,
    escapeinside={(*@}{@*)}, % Allows inline LaTeX
}


\newcommand{\myparagraph}[1]{\vspace{0.5ex}\noindent \textbf{#1. }}

%% Placeholder for numbers that depend on experiments not yet run.
%% Grep for \PENDING before submitting -- none must remain.
\newcommand{\PENDING}[1]{\textcolor{red}{\textbf{[PENDING: #1]}}}
